\begin{ouverture}[statistiques]{10}{Statistiques et diagrammes}{Partie 3 · Données et hasard}
{Effectifs · Fréquences · Moyenne · Médiane}
\objectif{Organiser et représenter des données.}
\objectif{Calculer une médiane sur une petite série.}
\objectif{Interpréter moyennes, fréquences et diagrammes circulaires.}
\prerequis{Pourcentages et angles.}
\end{ouverture}
\pagestyle{otMain}\essentiel\fontsize{11}{15}\selectfont
\phantomsection\label{cours-10-1}
\essSection{1}{Décrire une série}
\begin{definition}L’effectif d’une valeur est le nombre de fois où elle apparaît. L’effectif total est la somme des effectifs. La fréquence d’une valeur est son effectif divisé par l’effectif total. Les fréquences de toutes les catégories distinctes ont pour somme 1.\end{definition}
\begin{propriete}La moyenne se calcule en divisant la somme de toutes les valeurs par leur nombre. Si les valeurs ont des effectifs, on calcule la somme des produits valeur $\times$ effectif puis on divise par l’effectif total.\end{propriete}
\begin{definition}Après avoir rangé les valeurs, une médiane partage la série en deux parties contenant chacune au moins la moitié des données, au sens « inférieures ou égales » et « supérieures ou égales ». Pour un effectif impair, on prend la valeur centrale ; pour un effectif pair, ce livre prend la moyenne des deux valeurs centrales.\end{definition}
\begin{exemple}Pour $2,4,4,7,8$, la médiane est la 3e valeur, soit 4 ; la moyenne vaut $25/5=5$. Pour $2,4,7,9$, on choisit comme médiane $(4+7)/2=5{,}5$. La médiane n’est pas toujours une valeur observée.\end{exemple}
\clearpage
\phantomsection\label{cours-10-2}
\essSection{2}{Représenter et interpréter}
\begin{propriete}Dans un diagramme circulaire complet, l’angle d’un secteur vaut fréquence $\times360^\circ$. La somme des angles vaut $360^\circ$. Une fréquence de 25\,\% donne ainsi $90^\circ$.\end{propriete}
\label{fig-cours-10-2-cours-camembert}\figureGeo{cours-camembert}

\begin{savoirfaire}[Choisir un graphique]Un diagramme en bâtons compare des effectifs par valeur. Un diagramme circulaire montre le partage d’un total en catégories. Ajouter titre, légende, unités et graduations. Un tableur-grapheur peut construire ces représentations à partir d’un tableau.\end{savoirfaire}
\begin{attention}Toujours identifier le total de référence. Deux groupes de tailles différentes ne se comparent pas seulement avec leurs effectifs. La moyenne est sensible aux valeurs très grandes ; la médiane renseigne sur le centre de la série ordonnée, sans dire que toutes les valeurs sont proches.\end{attention}

\clearpage\rubrique{Découverte}
\begin{blocLivre}{ex-10-1}\begin{exercice}[1]{Un effectif total}
Les notes 8, 10, 12 ont respectivement pour effectifs 3, 5, 2. Donner l’effectif total et la fréquence de la note 10.
\end{exercice}\end{blocLivre}
\begin{blocLivre}{ex-10-2}\begin{exercice}[1]{Une moyenne simple}
Calculer la moyenne de 6, 9, 10 et 15.
\end{exercice}\end{blocLivre}
\begin{blocLivre}{ex-10-3}\begin{exercice}[1]{Une médiane impaire}
Déterminer une médiane de la série : 12, 7, 9, 15, 8.
\end{exercice}\end{blocLivre}
\begin{blocLivre}{ex-10-4}\begin{exercice}[1]{Une médiane paire}
Déterminer la médiane choisie dans ce livre pour : 3, 4, 8, 10, 12, 15.
\end{exercice}\end{blocLivre}
\begin{blocLivre}{ex-10-5}\begin{exercice}[1]{Un secteur circulaire}
Un sport regroupe 30\,\% des inscrits. Quel angle lui attribuer dans un diagramme circulaire complet ?
\end{exercice}\end{blocLivre}

\rubrique{Entraînement}
\begin{blocLivre}{ex-10-6}\begin{exercice}[2]{Une moyenne avec effectifs}
Les valeurs 2, 5 et 8 ont pour effectifs 4, 3 et 3. Calculer la moyenne.
\end{exercice}\end{blocLivre}
\begin{blocLivre}{ex-10-7}\begin{exercice}[2]{Lire un diagramme en bâtons}
Le diagramme donne les nombres de livres lus par 12 élèves. Déterminer la médiane.
\label{fig-ex-10-7}\figureGeo{stat-livres}
\end{exercice}\end{blocLivre}
\begin{blocLivre}{ex-10-8}\begin{exercice}[2]{Construire le partage}
Sur 40 élèves, 20 choisissent le football, 12 le basket et 8 le tennis. Construire le diagramme circulaire.
\label{fig-ex-10-8}\figureGeo{stat-sports}
\end{exercice}\end{blocLivre}
\begin{blocLivre}{ex-10-9}\begin{exercice}[2]{Lire un pourcentage}
Dans le diagramme, le secteur A mesure $144^\circ$. Il y a 75 réponses au total. Combien correspondent à A ?
\label{fig-ex-10-9}\figureGeo{stat-secteur-a}
\end{exercice}\end{blocLivre}
\begin{blocLivre}{ex-10-10}\begin{exercice}[2]{Une formule de tableur}
Les valeurs sont en A2:A4 et leurs effectifs en B2:B4. Proposer une méthode avec une colonne C pour calculer la moyenne.
\end{exercice}\end{blocLivre}

\rubrique{Problèmes}
\begin{blocLivre}{ex-10-11}\begin{exercice}[3]{Ajouter une valeur}
Cinq élèves ont une moyenne de 12 à un jeu. Un sixième rejoint le groupe avec 18. Quelle est la nouvelle moyenne ?
\end{exercice}\end{blocLivre}
\begin{blocLivre}{ex-10-12}\begin{exercice}[3]{Une valeur manquante}
Les nombres 8, 11, 14 et $x$ ont une moyenne de 12. Trouver $x$, puis la médiane de la série.
\end{exercice}\end{blocLivre}
\begin{blocLivre}{ex-10-13}\begin{exercice}[3]{Une valeur extrême}
Comparer moyenne et médiane des séries A : 4, 5, 6, 7, 8 et B : 4, 5, 6, 7, 38.
\end{exercice}\end{blocLivre}
\begin{blocLivre}{ex-10-14}\begin{exercice}[3]{Comparer deux groupes}
Dans A, 18 élèves sur 24 pratiquent un sport. Dans B, 21 sur 30 en pratiquent un. Quel groupe a la plus forte proportion de sportifs ?
\end{exercice}\end{blocLivre}
\begin{blocLivre}{ex-10-15}\begin{exercice}[3]{Interpréter une médiane}
Dans une série de 20 durées, la médiane choisie vaut 15 min. Peut-on conclure que dix durées valent exactement 15 min ? Donner un exemple qui justifie.
\end{exercice}\end{blocLivre}

\clearpage
